\documentclass[11pt]{article}

\usepackage[margin=1in]{geometry}
\usepackage[T1]{fontenc}
\usepackage{lmodern}
\usepackage{amsmath,amssymb}
\usepackage[colorlinks=true,linkcolor=blue,citecolor=blue,urlcolor=blue]{hyperref}

\emergencystretch=3em

\title{Literature Status: Negative Answer to Koldobsky's Measure-Slicing Problem}
\author{Run \texttt{fa\_banach\_001}, lane 8}
\date{2026-07-19}

\begin{document}
\maketitle

\noindent\textbf{Status.}
\texttt{literature\_already\_answered (negative; fails already for \(k=1\))}.
This note records an already-known later-literature answer, not a new proof
from the run.

\section*{Source Question}

The source paper is A. Koldobsky,
\emph{Slicing inequalities for measures of convex bodies}, arXiv:1412.8550,
published in \emph{Advances in Mathematics} 283 (2015), 473--488.
In Problem \texttt{prob} in the Introduction, Koldobsky asks whether there is
an absolute constant \(C\) such that, for every \(n\), every \(1\leq k<n\),
every origin-symmetric convex body \(L\subset\mathbb R^n\), and every measure
\(\mu\) with non-negative even continuous density,
\[
  \mu(L)
  \leq
  C^k
  \max_{H\in Gr_{n-k}} \mu(L\cap H)\, |L|^{k/n}.
\]
The same question is stated in the abstract of the source paper.

\section*{Supporting Answer}

The supporting paper is B. Klartag and G. V. Livshyts,
\emph{The lower bound for Koldobsky's slicing inequality via random rounding},
arXiv:1810.06189.
The title and introduction explicitly identify the paper as studying the
lower bound for Koldobsky's slicing inequality, and the bibliography cites
Koldobsky's 2015 paper.

For a measure \(\mu\) with continuous density and an origin-symmetric convex
body \(K\), Klartag--Livshyts define
\[
  S_{\mu,K}
  =
  \inf_{\xi\in S^{n-1}}
  \frac{\mu(K)}
       {|K|^{1/n}\mu^+(K\cap \xi^\perp)}.
\]
They recall that Koldobsky proved the general upper bound \(S_n\leq\sqrt n\).
Their Theorem 1 proves the matching lower bound: for every \(n\), there exist
a measure \(\mu\) and a symmetric convex body \(L\subset \mathbb R^n\) such
that for all \(\xi\in S^{n-1}\) and all \(t\geq 0\),
\[
  \mu^+(L\cap(\xi^\perp+t\xi))
  \leq
  \frac{C}{\sqrt n}\,\mu(L)|L|^{-1/n},
\]
where \(C\) is universal.

\section*{Identification}

Take \(t=0\) in the supporting theorem. Then every central hyperplane section
satisfies
\[
  \mu^+(L\cap \xi^\perp)
  \leq
  \frac{C}{\sqrt n}\,\mu(L)|L|^{-1/n}.
\]
Equivalently, \(S_{\mu,L}\geq c\sqrt n\). If the source problem had a
dimension-free constant for all \(k\), then the case \(k=1\) would imply
\[
  \mu(L)
  \leq
  C_0 \max_{\xi\in S^{n-1}}
  \mu^+(L\cap \xi^\perp)|L|^{1/n}
  \leq
  \frac{C_0 C}{\sqrt n}\,\mu(L),
\]
which is impossible for large \(n\). Thus the proposed absolute-constant
inequality in Problem \texttt{prob} is false already in codimension one.

The density condition matches the source problem. In the proof,
Klartag--Livshyts construct the measure as a convolution of the standard
Gaussian measure with an explicitly symmetrized discrete measure, so the
resulting density is continuous, non-negative, and even.

\section*{Scope}

This packet answers only the arbitrary-measure absolute-constant question in
arXiv:1412.8550. It does not affect Koldobsky's positive results for
unconditional bodies, for duals of bodies with bounded volume ratio, or for
proportional codimension \(k\geq\lambda n\). It also does not address
Bourgain's original Lebesgue-volume slicing problem.

For \(k=1\), the later paper gives the sharp order up to constants:
Koldobsky's upper bound gives \(S_n\leq\sqrt n\), and Klartag--Livshyts give
\(S_n\geq c\sqrt n\).

\section*{Evidence}

The run indexes were searched for arXiv:1412.8550, the source title, and core
terms including \texttt{slicing inequalities}, \texttt{measures of convex
bodies}, \texttt{Koldobsky's slicing inequality}, and \texttt{random rounding}.
No exact prior packet for arXiv:1412.8550 was found. Related prior slicing
packets concern Bourgain's volume slicing problem and do not duplicate this
measure-slicing status.

The source PDF is included as \texttt{source\_paper.pdf}. The supporting PDF is
included as \texttt{supporting\_paper\_1810.06189.pdf}.

\begin{thebibliography}{9}
\bibitem{Koldobsky2015}
A. Koldobsky,
\emph{Slicing inequalities for measures of convex bodies},
Advances in Mathematics 283 (2015), 473--488; arXiv:1412.8550.

\bibitem{KlartagLivshyts}
B. Klartag and G. V. Livshyts,
\emph{The lower bound for Koldobsky's slicing inequality via random rounding},
arXiv:1810.06189.
\end{thebibliography}

\end{document}

